\documentclass[border=4pt]{standalone}
% Shared extension of the installed skill's publication-grade preamble.
% TraceBank figure setup: packages, colour values and plot defaults.
\usepackage{tikz,pgfplots,amsmath,siunitx}
\pgfplotsset{compat=1.18}
\usetikzlibrary{positioning,calc,arrows.meta,fit,backgrounds,shapes.geometric}
\definecolor{cbBlue}{RGB}{0,114,178}
\definecolor{cbOrange}{RGB}{230,159,0}
\definecolor{cbGreen}{RGB}{0,158,115}
\definecolor{cbVermillion}{RGB}{213,94,0}
\definecolor{cbPurple}{RGB}{204,121,167}
\pgfplotsset{every axis/.append style={font=\small,line width=.6pt,
 tick style={black},grid=major,grid style={gray!25,line width=.3pt},
 legend style={font=\footnotesize,draw=gray!50,fill=white},
 label style={font=\small}}}

\usepackage{fontspec}
\usepackage{xeCJK}
\setmainfont{Arial}
\setCJKmainfont{Microsoft YaHei}
\setsansfont{Arial}
\setCJKsansfont{Microsoft YaHei}
\definecolor{ink}{HTML}{17343E}
\definecolor{muted}{HTML}{6C7C80}
\tikzset{
  every node/.style={font=\fontsize{8.3}{12}\selectfont,text=ink},
  box/.style={draw=cbBlue!45,fill=cbBlue!5,rounded corners=2mm,
    minimum height=11mm,text width=31mm,align=center,inner sep=3mm,line width=.6pt},
  safe/.style={box,draw=cbGreen!65!black,fill=cbGreen!7},
  caution/.style={box,draw=cbOrange!80!black,fill=cbOrange!9},
  strong/.style={box,draw=cbVermillion!80!black,fill=cbVermillion!7},
  neutral/.style={box,draw=muted!45,fill=muted!4},
  arrow/.style={-{Stealth[length=2mm,width=1.5mm]},line width=.65pt,draw=ink!75},
  response/.style={arrow,dashed,draw=cbBlue!70},
  note/.style={align=left,font=\fontsize{7.5}{11}\selectfont,text=muted},
  tinylabel/.style={font=\fontsize{7.2}{9}\selectfont,text=muted,fill=white,inner sep=1mm},
  title/.style={anchor=west,font=\fontsize{13}{17}\selectfont\bfseries,text=ink},
  subtitle/.style={anchor=west,font=\fontsize{8}{11}\selectfont,text=muted},
  boundary/.style={draw=cbBlue!45,dashed,rounded corners=3mm,line width=.65pt}
}

% ===== TUNABLES =====
\def\figwidth{183mm}
\def\diagramtitle{F4 · 任务依赖图与持久化节点状态}
% Source: validate_nodes / Orchestrator.run / control / release_handoff
\begin{document}
\begin{tikzpicture}[x=10mm,y=10mm]
\useasboundingbox (0,0) rectangle (18.3,11.4);
\node[title] at (0,11.1) {\diagramtitle};
\node[subtitle] at (0,10.5) {合法候选计划示例；当前编排器顺序执行就绪节点，不宣称并发调度};
\node[neutral,text width=65mm] (goal) at (9.1,9.1) {用户目标：先看账单，再给李娜转 300 元并查看卡片\\校验工具白名单、严格参数、依赖引用与无环性};
\node[safe,text width=51mm] (bills) at (9.1,7.1) {节点 A · bills\\只读账单分类与异常证据};
\node[caution,text width=51mm] (transfer) at (4.65,4.7) {节点 B · transfer\\depends\_on: [A]\\进入确认或验证状态后等待};
\node[safe,text width=51mm] (cards) at (13.55,4.7) {节点 C · cards\\depends\_on: [A]\\查询当前卡片状态};
\node[neutral,text width=65mm] (end) at (9.1,2.1) {聚合任务状态与业务回执\\这里只汇总终态，不是额外可执行工具};
\draw[arrow] (goal) -- (bills);
\draw[arrow] (bills.south) -- ++(0,-.35) -| (transfer.north);
\draw[arrow] (bills.south) -- ++(0,-.35) -| (cards.north);
\draw[arrow] (transfer.south) -- ++(0,-.55) -| (end.north);
\draw[arrow] (cards.south) -- ++(0,-.55) -| (end.north);
\node[note,text width=38mm,anchor=north west] at (.4,8.6) {图示表达依赖能力。离线解析器保留显式先后顺序；独立分支可由经校验的候选 DAG 表达。};
\node[note,text width=38mm,anchor=north west] at (14.15,8.6) {A 失败或取消时，依赖它的未执行后续步骤停止。已成功的资金操作不能用改状态来撤销。};
\node[note,text width=176mm,anchor=north west] at (.15,.95) {节点可持久化为 pending、awaiting\_confirmation、awaiting\_verification、unknown、succeeded、failed 或 cancelled。暂停与人工接管控制整个计划；操作员释放接管后仍需客户授权。重试和恢复属于节点状态管理，不在依赖图中画回边。};
\end{tikzpicture}
\end{document}
